\documentclass[11pt,a4paper]{article}
\usepackage[utf8]{inputenc}
\usepackage[T1]{fontenc}
\usepackage{lmodern}
\usepackage{amsmath,amssymb,amsfonts}
\usepackage{booktabs}
\usepackage{geometry}
\usepackage{hyperref}

\geometry{margin=2.2cm}

\title{\textbf{Topological Emergence and Dimensional Projection of Time:\\
An Effective Variational Framework}}
\author{E. R. dos Santos\\
\small Independent Research $\cdot$ Version 1.2.2 $\cdot$ License: GNU GPLv3}
\date{November 2026}

\begin{document}

\maketitle

\begin{abstract}
We present an exploratory effective field framework investigating the emergence of spatial metric scales, gauge interactions, gravitational rigidity, and Standard Model mass hierarchies from discrete topological defects in a primordial temporal substrate. Rather than treating matter, gauge fields, and spacetime as independent entities, localized particles are modeled as topologically confined temporal folds, gauge forces as geometric phase connections between adjacent lattice cells, and gravitation as the macroscopic elastic response of the emergent medium. Starting from the Heisenberg quantum non-nullity bound ($S_0 = 1/2$) and Kolmogorov--Arnold--Moser (KAM) stability of the golden ratio ($\phi$), crystallization into dense 3D ($z_{3d} = 12$) and 4D ($z_{4d} = 24$) packings induces a variational curvature defect $d^* = \frac{1}{2}\phi^{-1/3}$ and an Euler characteristic shift $\chi = 2.1$. Exponential dimensional projection from the Planck scale ($t_{\text{Planck}} \to \tau_c$) establishes the hadronic infrared scale $\mu_{\text{IR}} \approx 234.97\text{ MeV}$. Expanding the unified effective Lagrangian yields Leading-Order (LO) and Next-to-Leading-Order (NLO) spinor double-cover recoil predictions across ten Standard Model benchmarks, three quantum-atomic observables, and two falsifiable beyond-Standard-Model (BSM) states.
\end{abstract}

\section{Conceptual Postulates and Pre-Geometric Substrate}

In alignment with induced-gravity and geometrodynamic paradigms~\cite{Sakharov1967}, this framework treats temporal phase evolution as the fundamental pre-geometric substrate from which spatial dimensions and physical interactions emerge via spontaneous symmetry breaking:

\begin{enumerate}
    \item \textbf{Quantum Non-Nullity and Spontaneous Symmetry Breaking:} A strictly static, zero-action vacuum state ($\Delta E = 0, \Delta t = 0$) is precluded by the Heisenberg uncertainty bound $\Delta E \cdot \Delta t \ge \hbar/2$. Consequently, the trivial state $S = 0$ is quantum-mechanically unstable, forcing a primordial phase rupture bounded below by the minimal action $S_0 = 1/2$.
    
    \item \textbf{KAM Selection of the Golden Ratio:} In the resulting non-linear pre-spatial phase dynamics, rational frequency modes undergo resonant destabilization. By the Kolmogorov--Arnold--Moser (KAM) theorem~\cite{Kolmogorov1954,Moser1962,Arnold1963}, the invariant torus associated with the maximally irrational winding number—the golden ratio $\phi = (1 + \sqrt{5})/2$—survives as the unique dynamically stable attractor.
    
    \item \textbf{Topological Folding and Lattice Crystallization:} Phase folding across the temporal substrate creates localized negative-curvature defects that require topological accommodation by the surrounding medium. At the critical percolation threshold, maximal sphere-packing coordination emerges in 3D icosahedral/FCC ($z_{3d} = 12$) and 4D $D_4$ ($z_{4d} = 24$) lattice geometries. Because icosahedral symmetry cannot tile a spherical topology ($\chi_0 = 2$) without frustration, a topological vortex density arises:
    \begin{equation}
        \rho_v^* = \frac{1}{z_{3d} - \chi_0} = \frac{1}{10} = 0.1, \qquad \chi = \chi_0 + \rho_v^* = 2.1.
    \end{equation}
    
    \item \textbf{Dimensional Projection and Induced Gravitational Rigidity:} The bare geometric gauge coupling $\alpha_{\text{bare}}^{-1}$ governs an exponential scaling from the basal Planck time $t_{\text{Planck}} = \sqrt{\hbar G / c^5}$ to the coherent hadronic timescale $\tau_c$, projecting onto the spatial infrared scale $\mu_{\text{IR}}$:
    \begin{equation}
        \alpha_{\text{bare}}^{-1} = \frac{2\pi\chi\sqrt{z_{4d}}}{d^*} \left( 1 - \frac{d^*}{\sqrt{z_{4d}}} - \frac{\chi}{z_{3d}^2} \right), \qquad \tau_c = t_{\text{Planck}} \exp\left( \frac{\alpha_{\text{bare}}^{-1} - \chi/z_{3d}}{3} \right), \qquad \mu_{\text{IR}} = \frac{\hbar c}{c\tau_c}.
    \end{equation}
    Here, fundamental constants ($\hbar, c, G$) serve strictly as Planck-scale unit conversion parameters mapping dimensionless topological ratios into physical energy units (GeV/MeV). In this unified ontology, baryonic and fermionic matter correspond to topologically confined temporal folds, gauge interactions represent geometric phase connections across coordination neighborhoods, and gravitation manifests as the macroscopic elastic rigidity of the temporal mesh reacting to localized concentrations of confined energy~\cite{Sakharov1967}.
    
    \item \textbf{Spinor Double-Cover Topology:} Spin-1/2 fermions correspond to negative-curvature topological folds requiring a $4\pi$ double-cover rotation across the $2z_{4d} = 48$ local spinorial states of the $D_4$ lattice.
\end{enumerate}

\section{Unified Effective Lagrangian and Variational Derivations}

The continuum limit of the emergent lattice is governed by the effective action $S = \int d^4x \sqrt{-g_{\text{eff}}}\,\mathcal{L}_{\text{Total}}$ composed of four coupled sectors:
\begin{equation}
    \mathcal{L}_{\text{Total}} = \mathcal{L}_{\text{Lattice}}(\Phi, d^*) + \mathcal{L}_{\text{Gauge}}(F_{\mu\nu}, \chi) + \mathcal{L}_{\text{Spinor}}(\psi, \Phi, d^*) + \mathcal{L}_{\text{Topological}}(a, G\tilde{G}).
\end{equation}

\subsection{Variational Origin of the Discrete Curvature Defect ($d^*$)}
Let $x = d / S_0 = 2d$ denote the dimensionless curvature defect normalized by the minimal quantum action. Across $D = 3$ emergent isotropic spatial dimensions, the spatial defect volume scales as $x^3$, whereas self-interaction across the 6D phase space (position and conjugate momentum) scales as $x^6$. Extremizing the geometric frustration functional $\mathcal{F}(x, \Lambda) = \frac{1}{2}x^2 - \Lambda(x^3 + x^6 - 1)$ with respect to the Lagrange multiplier $\Lambda$ ($\delta\mathcal{F} = 0$) enforces the unitary phase-space conservation constraint $x^3 + x^6 = 1$. Setting $x^3 = \phi^{-1}$ identically satisfies the golden ratio identity $\phi^{-1} + \phi^{-2} = 1$, fixing the unique positive stationary point at $x^* = \phi^{-1/3}$:
\begin{equation}
    d^* = S_0 \phi^{-1/3} = \frac{1}{2}\phi^{-1/3} \approx 0.425900.
\end{equation}

\subsection{Lattice Symmetry Breaking and Gauge-Hadronic Sectors}
The scalar order parameter $\Phi$ represents the collective elastic mode of the lattice, acquiring a vacuum expectation value $v_{\text{Higgs}} = f_\pi \alpha^{-1}(0) z_{4d} \sqrt{\chi}\,\eta^{3/2}$, where $f_\pi = \mu_{\text{IR}}/\sqrt{z_{4d}}$ and $\eta = (1 - d^*)^{-1/z_{3d}}$ accounts for vacuum phase screening across the $z_{3d}$ bonds:
\begin{equation}
    \mathcal{L}_{\text{Lattice}} = \frac{1}{2}\left(1 - \frac{d^*}{\sqrt{z_{4d}}} - \frac{\chi}{z_{3d}^2}\right)(D_\mu \Phi)^\dagger (D^\mu \Phi) - \lambda_H \left(\Phi^\dagger \Phi - \frac{v_{\text{Higgs}}^2}{2}\right)^2.
\end{equation}
Projecting a four-point scalar vertex from the $z_{4d}$ connections onto the $z_{3d}$ spatial links and minimizing the effective potential ($\partial V_{\text{eff}}/\partial(\Phi^\dagger\Phi) = 0$) fixes the quartic coupling at $\lambda_H = \frac{1}{4}(z_{3d}/z_{4d})\eta = \frac{z_{3d}}{4z_{4d}}\eta$.

In the gauge sector, a central lattice site plus its immediate 3D neighborhood spans $N_{\text{local}} = z_{3d} + 1 = 13$ orthogonal local states. During electroweak symmetry breaking $SU(2)_L \times U(1)_Y \to U(1)_{\text{em}}$, aligning the three weak-isospin generators with the $D = 3$ emergent spatial axes fixes the spectral projection ratio:
\begin{equation}
    \sin^2\theta_W = \frac{D}{z_{3d} + 1} = \frac{3}{13} \approx 0.23077.
\end{equation}
Simultaneously, baryonic matter is stabilized as topological Skyrme solitons~\cite{Skyrme1961} in the chiral-gauge Lagrangian, with Skyrme coupling $e_{\text{Skyrme}} = (d^*)^{-1/4}$ and topological soliton invariant $c_{\text{sol}} = 3^5/10 = 24.3$, yielding $M_p = (c_{\text{sol}}/e_{\text{Skyrme}})f_\pi$:
\begin{equation}
    \mathcal{L}_{\text{Gauge}} = -\frac{1}{4g_{\text{bare}}^2} F^a_{\mu\nu} F^{a\,\mu\nu} + \frac{f_\pi^2}{4}\text{Tr}(\partial_\mu U \partial^\mu U^\dagger) + \frac{1}{32 e_{\text{Skyrme}}^2}\text{Tr}[U^\dagger \partial_\mu U, U^\dagger \partial_\nu U]^2.
\end{equation}

\subsection{Spinor Kinetic Recoil: Perturbative Origin of LO and NLO}
Propagation of a spin-1/2 topological fold $\psi$ across the $2z_{4d} = 48$ double-cover states induces a first-order curvature modification in the Dirac kinetic operator:
\begin{equation}
    \mathcal{L}_{\text{Spinor}} = \bar{\psi} \left[ i\gamma^\mu D_\mu \left( 1 + \frac{d^*}{2z_{4d}} + \mathcal{O}(d^{*2}) \right) - y_f(\eta, \chi)\Phi \right] \psi.
\end{equation}
Truncating at zeroth order defines the Leading-Order (LO) spectrum. Retaining the first-order kinetic renormalization $Z_\psi = 1 + d^*/(2z_{4d})$ and canonically rescaling $\psi_{\text{can}} = Z_\psi^{1/2}\psi$ transforms the Yukawa mass term as $m_f^{\text{NLO}} = m_f^{\text{LO}} Z_\psi^{-1} \approx m_f^{\text{LO}}(1 - \frac{d^*}{2z_{4d}})$, defining the Next-to-Leading-Order (NLO) charged-lepton masses. One-loop QED vacuum polarization is evaluated self-consistently up to $\Lambda = m_\tau$ (determined via the geometric Koide condition $Q = 2/3$~\cite{Koide1982}), while the top quark saturates the emergent vacuum via $m_t = \frac{v_{\text{Higgs}}}{\sqrt{2}}(1 - d^*/z_{3d}^2)$, entering $M_W$ and $M_Z$ through standard $\Delta\rho$ radiative corrections.

\section{Numerical Results and Empirical Comparison}

Table~\ref{tab:sm_benchmarks} compares the LO and NLO evaluations against Particle Data Group (PDG) experimental values~\cite{PDG2024}. Theoretical uncertainties ($\sigma_{\text{tot}}$) combine renormalization scale variation ($\mu \in [0.5, 2.0]\,m_\tau$), higher-order radiative truncation ($\alpha/\pi$), and experimental uncertainty in $G$ ($\Delta G/G = 2.2 \times 10^{-5}$).

\begin{table}[htbp]
\centering
\caption{Comparison of Leading-Order (LO) and Next-to-Leading-Order (NLO) model outputs against experimental PDG reference values~\cite{PDG2024}.}
\label{tab:sm_benchmarks}
\begin{tabular}{lcccc}
\toprule
\textbf{Observable} & \textbf{LO (Tree-Level)} & \textbf{NLO ($\pm\sigma_{\text{tot}}$)} & \textbf{PDG Reference} & \textbf{NLO Tension} \\
\midrule
$M_p$ (GeV) & 0.9416 & $0.9416 \pm 0.0022$ & 0.9383 & $1.5\sigma$ \\
$\alpha^{-1}(0)$ & 137.1069 & $137.1050 \pm 0.2943$ & 137.0360 & $0.2\sigma$ \\
$m_e$ (MeV) & 0.5182 & $0.5136 \pm 0.0016$ & 0.5110 & $1.6\sigma$ \\
$m_\mu$ (MeV) & 106.5782 & $105.6326 \pm 0.2452$ & 105.6584 & $0.1\sigma$ \\
$m_\tau$ (MeV) & 1793.2838 & $1777.3743 \pm 4.1399$ & 1776.8600 & $0.1\sigma$ \\
$v_{\text{Higgs}}$ (GeV) & 245.1447 & $245.1413 \pm 0.7752$ & 246.2197 & $1.4\sigma$ \\
$m_t$ (GeV) & 172.8308 & $172.8284 \pm 0.5465$ & 172.6900 & $0.3\sigma$ \\
$M_W$ (GeV) & 80.6378 & $80.6373 \pm 0.2286$ & 80.3770 & $1.1\sigma$ \\
$M_Z$ (GeV) & 91.2262 & $91.2256 \pm 0.2586$ & 91.1876 & $0.1\sigma$ \\
$M_H$ (GeV) & 125.4396 & $125.4379 \pm 0.3966$ & 125.2500 & $0.5\sigma$ \\
\bottomrule
\end{tabular}
\end{table}

\subsection{Emergent Quantum Scales and Testable BSM Benchmarks}
Without external QED loop coefficients, geometric defect recoil yields an electron anomalous magnetic moment $a_e = \frac{\alpha}{2\pi} - d^*\left(\frac{\alpha}{\pi}\right)^2 = 0.0011585$ ($0.10\%$ deviation from $0.0011597$; $0.3\sigma$), Bohr radius $a_0 = 0.5267\text{ \AA}$ ($0.46\%$ vs.\ $0.5292\text{ \AA}$; $1.0\sigma$), and Rydberg energy $R_\infty = 13.6621\text{ eV}$ ($0.41\%$ vs.\ $13.6057\text{ eV}$; $0.6\sigma$). Beyond the Standard Model, the lattice predicts:
\begin{itemize}
    \item \textbf{Heavy Icosahedral Scalar ($M_{H_2}$):} $M_{H_2} = v_{\text{Higgs}}\sqrt{z_{3d}/2} \approx 600.47 \pm 1.90\text{ GeV}$, accessible to heavy Higgs resonance searches at the HL-LHC~\cite{HLLHC2019}.
    \item \textbf{Topological QCD Axion ($m_a$):} At the classical post-inflationary Peccei--Quinn scale $f_a = v_{\text{Higgs}}(\alpha^{-1}(0))^3 z_{4d} \approx 1.52 \times 10^{10}\text{ GeV}$, the topological sector yields a QCD axion of mass $m_a \approx 0.376\text{ meV}$ ($375.8\,\mu\text{eV}$, corresponding to $\nu \approx 90.9\text{ GHz}$), squarely within the target mass window ($40\text{--}400\,\mu\text{eV}$) of dielectric haloscopes such as MADMAX~\cite{MADMAX2017}.
\end{itemize}

\section{Methodological Limitations and Transparency}
Because the effective algebraic operators in this phenomenological model were constructed with prior knowledge of experimental benchmarks (\emph{post-diction}), this work does not claim immunity to combinatorial selection effects (\emph{look-elsewhere effect}). Complete physical validation depends on: (i) deriving the discrete $D_4$ lattice crystallization dynamically from a fundamental pre-geometric action, and (ii) experimental confirmation or refutation of the predicted candidate states ($M_{H_2} \approx 600.47\text{ GeV}$ and $m_a \approx 0.376\text{ meV}$).

\medskip
\noindent\textbf{Transparency and Methodology Note:} The conceptual formulation, theoretical axioms, and architecture of this model were idealized by E. R. dos Santos. The logical implementation in Python, mathematical structuring, and code refactoring were developed with the assistance of Artificial Intelligence. This work does not aim, in any way or at any time, to exhaust the subject, supersede similar research, or provide definitive answers of any kind under any aspect. All numerical results, equations, and computational routines are available for reading, auditing, testing, and modification under the terms of the GNU General Public License v3.0 (GPLv3).

\begin{thebibliography}{99}

\bibitem{Kolmogorov1954}
A.~N.~Kolmogorov,
``On conservation of conditionally periodic motions for a small change in Hamilton's function,''
\textit{Dokl. Akad. Nauk SSSR} \textbf{98}, 527--530 (1954).

\bibitem{Moser1962}
J.~Moser,
``On invariant curves of area-preserving mappings of an annulus,''
\textit{Nachr. Akad. Wiss. G\"ottingen Math.-Phys. Kl. II} \textbf{1962}, 1--20 (1962).

\bibitem{Arnold1963}
V.~I.~Arnold,
``Proof of a theorem of A.~N.~Kolmogorov on the invariance of quasi-periodic motions under small perturbations of the Hamiltonian,''
\textit{Russ. Math. Surv.} \textbf{18}(5), 9--36 (1963).

\bibitem{Sakharov1967}
A.~D.~Sakharov,
``Vacuum quantum fluctuations in curved space and the theory of gravitation,''
\textit{Dokl. Akad. Nauk SSSR} \textbf{177}, 70--71 (1967)
[\textit{Gen. Rel. Grav.} \textbf{32}, 365--367 (2000)].

\bibitem{Skyrme1961}
T.~H.~R.~Skyrme,
``A non-linear field theory,''
\textit{Proc. Roy. Soc. Lond. A} \textbf{260}, 127--138 (1961).

\bibitem{Koide1982}
Y.~Koide,
``Fermion-boson two-body model of quarks and leptons and Cabibbo mixing,''
\textit{Lett. Nuovo Cim.} \textbf{34}, 201--205 (1982).

\bibitem{PDG2024}
S.~Navas \textit{et al.} [Particle Data Group],
``Review of Particle Physics,''
\textit{Phys. Rev. D} \textbf{110}(3), 030001 (2024).

\bibitem{MADMAX2017}
A.~Caldwell \textit{et al.} [MADMAX Working Group],
``Dielectric Haloscopes: A New Way to Detect Axion Dark Matter,''
\textit{Phys. Rev. Lett.} \textbf{118}(9), 091801 (2017).

\bibitem{HLLHC2019}
M.~Cepeda \textit{et al.},
``Report from Working Group 2: Higgs Physics at the HL-LHC and HE-LHC,''
\textit{CERN Yellow Rep. Monogr.} \textbf{7}, 221--584 (2019).

\end{thebibliography}

\appendix
\section{Code Availability and Reproducibility}
The complete reference implementation in Python (\texttt{topological\_emergence.py}, v1.2.2) used to compute all Leading-Order (LO) and Next-to-Leading-Order (NLO) benchmarks, uncertainty propagations, and candidate beyond-Standard-Model observables is openly available for auditing, testing, and modification under the GNU General Public License v3.0 (GPLv3).

\end{document}
